% !TEX program = pdflatex
\documentclass[11pt]{article}

\usepackage[T1]{fontenc}
\usepackage[utf8]{inputenc}
\usepackage{lmodern}
\usepackage{microtype}
\usepackage[margin=1in]{geometry}
\usepackage{amsmath,amssymb,amsthm,mathtools}
\usepackage{booktabs}
\usepackage{tabularx}
\usepackage{array}
\usepackage{enumitem}
\usepackage{xcolor}
\usepackage[hidelinks]{hyperref}
\usepackage[nameinlink,noabbrev]{cleveref}

\newcolumntype{Y}{>{\raggedright\arraybackslash}X}
\newcolumntype{P}[1]{>{\raggedright\arraybackslash}p{#1}}

\newtheorem{theorem}{Theorem}[section]
\newtheorem{proposition}[theorem]{Proposition}
\newtheorem{lemma}[theorem]{Lemma}
\newtheorem{corollary}[theorem]{Corollary}
\theoremstyle{definition}
\newtheorem{definition}[theorem]{Definition}
\newtheorem{assumption}[theorem]{Assumption}
\theoremstyle{remark}
\newtheorem{remark}[theorem]{Remark}
\newtheorem{example}[theorem]{Example}

\DeclareMathOperator{\Tr}{Tr}
\DeclareMathOperator{\dist}{dist}
\DeclareMathOperator{\Ran}{Ran}
\DeclareMathOperator{\supp}{supp}
\newcommand{\id}{\mathrm{id}}
\newcommand{\C}{\mathbb{C}}
\newcommand{\R}{\mathbb{R}}
\newcommand{\N}{\mathbb{N}}
\newcommand{\B}{\mathcal{B}}
\newcommand{\T}{\mathcal{T}}
\newcommand{\W}{\mathsf{W}}
\newcommand{\axis}[1]{\mathsf{#1}}

\title{\textbf{Conditional Classification of Projection-Based Effective Theories:}\\
\large Witness Records, Comparison Axes, and the Limits of Inevitability in Modal Triplet Theory}
\author{Peter Nero}
\date{Version 3, July 2026}

\begin{document}
\maketitle

\begin{abstract}
Irreversible records, ultraviolet control, event-like outcomes, and consistent
matter sectors are genuine demands on physical theories. They do not, by
themselves, force all viable theories into one mathematical universality class.
The missing information is the mechanism: which states and operations are
admitted, what is projected, what counts as recovery, which dynamics stabilizes
records, how ultraviolet errors are controlled, how events are generated, and
which global gauge and bundle data define the matter sector. This paper replaces
an earlier inevitability claim by a conditional classification framework. A
theory is represented by a seven-axis witness record covering operational
measurement, descent and recovery, stability, ultraviolet control, event
structure, gauge and matter data, and empirical validation. We prove five
bounded results: noninjective descent has no left inverse; compact cores of open
record regions have positive margins; transverse crossings are locally finite;
representations of a specified compact \(U(1)\) have integer weights; and
multi-axis equivalence classes refine monotonically as comparison requirements
are added. Explicit counterexamples show why operational irreversibility need
not mean algebraic noninvertibility, low-energy predictivity does not select
smooth spectral damping, countable outcomes do not create a causal set, and
local anomaly cancellation or charge quantization does not determine the
observed Standard Model. Collapse models, monitored circuits, effective field
theory, asymptotic safety, loop quantum gravity, causal sets, noncommutative
geometry, and Modal Triplet Theory are therefore compared only through completed
witness rows. The result is a falsifiable taxonomy and release-ready comparison
protocol, not a proof that the apparent plurality of physical theories is
illusory.
\end{abstract}

\section*{Revision note for Version 3}

\begin{description}[leftmargin=3.3cm,style=nextline]
\item[Supersedes.] Version 2, DOI
\href{https://doi.org/10.5281/zenodo.18268193}{10.5281/zenodo.18268193}.
\item[Reason.] Version 2 presented five broad ``inevitability'' theorems and
concluded that all successful measurement, quantum-gravity, ultraviolet, and
matter frameworks belong to a single coherent universality class. Those
conclusions did not follow from the stated empirical requirements.
\item[Resolution.] Version 3 withdraws that classification, corrects the
direction of the inverse used in projection arguments, distinguishes
operational nonrecoverability from algebraic noninvertibility, replaces
automatic basins and thresholds by typed stability hypotheses, separates
renormalization from spectral damping, and restores the missing global gauge,
anomaly-inflow, charge-normalization, and bundle data.
\item[Retained content.] The paper retains a conditional witness-based taxonomy
and the correctly scoped local implications proved from completed witness rows.
\item[Remaining boundary.] No common ontology, ultraviolet completion, event
source, observed matter spectrum, gravity theory, or global physical equivalence
is inferred without a separate typed bridge and validation record.
\end{description}

This revision note records the release delta; it is not part of the abstract.

\tableofcontents

\section{The question that can actually be answered}
\label{sec:orientation}

Theories can resemble one another for several different reasons. Two microscopic
models may have the same infrared fixed point. Two measurement models may
produce the same outcome probabilities. Two compactifications may have the same
low-energy gauge algebra. Two operators may have nearby compressed resolvents.
None of these statements implies the others.

The word \emph{universality} is therefore incomplete unless it names:
\begin{enumerate}[label=(\arabic*)]
  \item the objects being compared;
  \item the maps that place them in one common space;
  \item the topology, norm, or operational metric of comparison;
  \item the parameter range and time horizon;
  \item the observables whose agreement matters; and
  \item the error and provenance of the comparison.
\end{enumerate}

This distinction is especially important in Modal Triplet Theory (MTT).
The current operator results establish controlled reduction inside a typed
model and local comparison after retained spaces have been explicitly
transported to a common Hilbert space
\cite{NeroCoherentReduction2026,NeroLocalRobustness2026}. They do not prove that
unrelated theories share one measurement dynamics, ultraviolet completion,
event ontology, gauge group, or matter content.

The aim here is modest but useful. We ask:
\begin{quote}
What finite record must be supplied before two proposed theories can be said to
belong to the same controlled class, and what follows from that record?
\end{quote}
This converts a broad philosophical slogan into an auditable classification
problem.

\section{Desiderata are not mechanisms}
\label{sec:desiderata}

\subsection{Five empirical demands}

For orientation, consider five demands often imposed on a candidate physical
theory:
\begin{itemize}[leftmargin=1.5em]
  \item \textbf{\(\axis{D}_{\rm rec}\): records.}
  It reproduces stable macroscopic records and their observed statistics.

  \item \textbf{\(\axis{D}_{\rm loc}\): locality.}
  It supplies the relevant relativistic locality or causal-compatibility
  property on the declared observables.

  \item \textbf{\(\axis{D}_{\rm pred}\): predictivity.}
  It gives finite, controlled predictions in a stated energy and accuracy
  regime.

  \item \textbf{\(\axis{D}_{\rm event}\): outcomes.}
  It explains or models individual registered outcomes when such outcomes are
  part of the intended ontology.

  \item \textbf{\(\axis{D}_{\rm matter}\): matter.}
  It reproduces the observed gauge representations, anomaly structure, charge
  assignments, masses, mixings, and forbidden or suppressed processes to its
  declared tier.
\end{itemize}

These demands are not yet mathematical premises. For example,
``irreversible record'' does not say whether irreversibility means a
noninjective map, absence of a completely positive recovery, thermodynamic cost,
or practical inaccessibility of environmental degrees of freedom. Likewise,
``finite predictivity'' does not choose among renormalization, an effective
cutoff, an ultraviolet fixed point, modular finiteness, a lattice continuum
limit, or nonlocal spectral damping.

\subsection{A claim ladder}

\begin{table}[htbp]
\centering
\small
\begin{tabularx}{\textwidth}{@{}P{0.20\textwidth}Y Y@{}}
\toprule
Level & What is established & What is still not established \\
\midrule
Vocabulary &
A common list of words such as projection, basin, gap, event, or topology. &
That the words denote the same mathematical objects. \\

Typed encoding &
Each theory supplies objects and maps with compatible domains and codomains. &
That the maps commute or preserve dynamics and observables. \\

Controlled comparison &
A norm, region, transport map, and quantitative error are proved. &
Global identity, arbitrary-time agreement, or physical equivalence. \\

Physical equivalence &
Preparations, dynamics, observables, units, uncertainties, and interpretation
are matched. &
Microscopic identity or unique ontology. \\

Selection &
One candidate is chosen by an independent source rule and passes held-out
tests. &
That every viable theory had to use the same mechanism. \\
\bottomrule
\end{tabularx}
\caption{A comparison must not be promoted by skipping a level.}
\label{tab:ladder}
\end{table}

Version 2 moved directly from shared vocabulary and empirical desiderata to
global physical classification. The remainder of this paper supplies the
missing intermediate record.

\section{The seven-axis witness record}
\label{sec:witness}

\begin{definition}[Classification witness]
\label{def:witness}
A \emph{classification witness} for a model \(a\) is a tuple
\[
 \W_a=
 \bigl(
   \axis{M}_a,\axis{P}_a,\axis{S}_a,\axis{U}_a,
   \axis{E}_a,\axis{G}_a,\axis{V}_a;\pi_a
 \bigr),
\]
with the following rows.
\begin{enumerate}[label=\(\axis{\Alph*}\),leftmargin=2.8em]
  \item \(\axis{M}_a\), the operational-measurement row: state space,
  admissible preparations, instruments, outcomes, record channel, and
  operational metric.

  \item \(\axis{P}_a\), the projection or descent row: source and effective
  spaces, descent map, declared recovery class, injectivity or recovery defect,
  and approximation error.

  \item \(\axis{S}_a\), the stability row: dynamics, domain, metric or norm,
  record regions or metastable sets, margins, retention estimate, and time
  horizon.

  \item \(\axis{U}_a\), the ultraviolet row: regulator or construction,
  renormalization and continuum prescription, symmetry identities, error or
  remainder, energy range, and status of ultraviolet completion.

  \item \(\axis{E}_a\), the event row: path or history space, event functional,
  stopping/reset rule, ordering relation, local-finiteness or dwell-time
  condition, and outcome law.

  \item \(\axis{G}_a\), the gauge-and-matter row: global gauge group,
  representations, charge normalization, spin and bundle data, anomaly
  functional including inflow or counterterms, and operator-selection rule.

  \item \(\axis{V}_a\), the validation row: preparation and observable maps,
  unit conventions, matching scale, uncertainty budget, comparison data, and
  held-out tests.
\end{enumerate}
The provenance row \(\pi_a\) contains source identifiers, versions, hashes,
proof status, and known open obligations.
\end{definition}

The rows deliberately overlap at their interfaces. An outcome law belongs to
the measurement row, while the event row says how individual events are
generated or encoded. A spectral gap belongs to the stability or operator
certificate, while ultraviolet control must explain what happens across scales.
The validation row prevents formal similarity from being mistaken for
empirical agreement.

\begin{definition}[Complete comparison on an axis set]
Let \(I\subseteq\{\axis{M},\axis{P},\axis{S},\axis{U},
\axis{E},\axis{G},\axis{V}\}\). A comparison of models \(a\) and \(b\) is
\emph{complete on \(I\)} when every row in \(I\) supplies:
\begin{enumerate}[label=(\roman*)]
  \item explicit transports to common comparison objects;
  \item a declared exact relation or quantitative tolerance;
  \item all hypotheses needed for that relation;
  \item an error or obstruction certificate; and
  \item source provenance.
\end{enumerate}
\end{definition}

Blank rows are not failures of a theory. They are failures of the proposed
comparison. This is the central bookkeeping correction.

\section{Projection, recovery, and operational irreversibility}
\label{sec:projection}

\subsection{Three notions that must not be conflated}

Let \(X\) be a source state space, \(Y\) an effective state space, and
\(P:X\to Y\) a descent map.
\begin{enumerate}[label=(\arabic*)]
  \item \(P\) is \emph{algebraically noninjective} if two source states have the
  same image.
  \item \(P\) is \emph{operationally nonrecoverable} relative to a class
  \(\mathfrak D\) if no allowed decoder in \(\mathfrak D\) recovers all source
  states.
  \item A process is \emph{thermodynamically or practically irreversible} if
  reversal lies outside a declared resource or control regime.
\end{enumerate}
Only the first is a property of the set map alone.

\begin{proposition}[Noninjective descent has no left inverse]
\label{prop:no-left-inverse}
If \(P:X\to Y\) is noninjective, there is no map \(L:Y\to X\) satisfying
\[
 L\circ P=\id_X.
\]
\end{proposition}

\begin{proof}
Choose \(x_1\neq x_2\) with \(P(x_1)=P(x_2)\). If a left inverse existed, then
\[
 x_1=L(P(x_1))=L(P(x_2))=x_2,
\]
a contradiction.
\end{proof}

\begin{remark}[Left versus right]
The obstructed object in \cref{prop:no-left-inverse} is a \emph{left} inverse
of \(P\). A right inverse \(R:Y\to X\), satisfying \(P\circ R=\id_Y\), is a
section and may exist for a surjective but noninjective map. Noninjectivity
therefore proves loss of unique recovery, not absence of every section.
\end{remark}

\subsection{Operational nonrecoverability can occur without algebraic loss}

Quantum operations are naturally described by completely positive maps and
instruments \cite{DaviesLewis1970}. A microscopic unitary dilation may coexist
with a reduced, operationally irreversible channel
\cite{Stinespring1955}. The following elementary example blocks the old
necessity argument.

\begin{example}[Invertible linear map with no physical recovery]
\label{ex:depolarizing}
On \(d\)-dimensional density matrices let
\[
 \Phi_p(\rho)=p\rho+(1-p)\frac{I_d}{d},
 \qquad 0<p<1.
\]
On the affine space of trace-one Hermitian matrices, \(\Phi_p\) has the linear
inverse
\[
 \Phi_p^{-1}(\sigma)
 =\frac{1}{p}\sigma-\frac{1-p}{pd}I_d.
\]
That inverse is not a quantum channel. Indeed, for orthogonal pure states
\(\rho,\sigma\), trace distance contracts from \(1\) to \(p\). If a completely
positive trace-preserving map \(D\) recovered both states, contractivity would
give
\[
 1=D_{\rm tr}(\rho,\sigma)
 =D_{\rm tr}\bigl(D\Phi_p(\rho),D\Phi_p(\sigma)\bigr)
 \leq D_{\rm tr}\bigl(\Phi_p(\rho),\Phi_p(\sigma)\bigr)
 =p,
\]
which is impossible.
\end{example}

Thus ``no admissible inverse'' is often the correct operational statement.
It depends on the recovery class. It is not equivalent to noninjectivity, and
neither statement follows from the bare observation that records look
irreversible. Objective-collapse models, for example, introduce explicit
modified dynamics rather than deriving it from a general projection theorem
\cite{GRW1986}. The current MTT measurement paper likewise treats measurement
as an ordinary physical interaction plus an explicitly typed completion and
record process \cite{NeroMeasurement2026}.

\begin{definition}[Recovery defect]
For a source set \(K\subseteq X\), metric \(d_X\), and allowed decoder class
\(\mathfrak D\), define
\[
 \delta_{\rm rec}(P;K,\mathfrak D)
 =
 \inf_{D\in\mathfrak D}\;
 \sup_{x\in K} d_X\bigl(D(Px),x\bigr).
\]
\end{definition}

This scalar is meaningful only after \(K\), \(d_X\), and \(\mathfrak D\) are
fixed. It distinguishes exact information loss from restricted or approximate
recovery.

\section{Record stability is not automatically a dynamical basin}
\label{sec:stability}

\subsection{Record regions and positive margins}

Let \((Y,d)\) be an effective state space and \(r:Y\to\mathcal R\) a record
label. For a label \(\alpha\), write \(B_\alpha=r^{-1}(\alpha)\).
If \(B_\alpha\) is open, nearby states share the same label. This is a
kinematic robustness statement. A dynamical basin additionally requires a
specified evolution and an attraction or retention property.

\begin{proposition}[Compact record cores have positive margins]
\label{prop:record-margin}
Let \(B_\alpha\) be open in a metric space \(Y\), and let
\(K_\alpha\subset B_\alpha\) be compact. Then
\[
 m_\alpha
 =
 \dist\bigl(K_\alpha,Y\setminus B_\alpha\bigr)>0,
\]
unless \(Y\setminus B_\alpha\) is empty, in which case the margin is
unbounded.
\end{proposition}

\begin{proof}
The complement \(Y\setminus B_\alpha\) is closed. The continuous function
\(y\mapsto \dist(y,Y\setminus B_\alpha)\) is strictly positive on
\(K_\alpha\). Compactness gives a positive minimum.
\end{proof}

The proposition explains exactly when a finite stability margin is justified.
It does not prove that every record region is open, that a useful compact core
exists, or that dynamics preserves it.

\begin{definition}[Finite-horizon record certificate]
A finite-horizon record certificate consists of a propagator or channel family
\(\{\Phi_t\}_{0\leq t\leq T}\), a record core \(K_\alpha\), a record region
\(B_\alpha\), a metric \(d\), and a bound
\[
 \inf_{y\in K_\alpha}
 \Pr_y\!\left[
   \Phi_t(y)\in B_\alpha\ \text{for all }0\leq t\leq T
 \right]
 \geq 1-\eta_\alpha.
\]
For deterministic dynamics the probability is omitted.
\end{definition}

Attraction, metastability, redundancy, superselection, and error correction can
all stabilize records, but they are different mechanisms. A sharp ``knee'' in
a control parameter is also not automatic. The smooth response
\[
 p(a)=\frac{1}{1+e^{-a}}
\]
changes between two record probabilities without a singular threshold. A phase
transition or first-exit boundary must be proved in the chosen model.
Measurement-induced entanglement transitions are genuine universality results
for specified monitored many-body models, not a theorem about all measurement
dynamics \cite{SkinnerRuhmanNahum2019}.

\section{Ultraviolet control is a family of mechanisms}
\label{sec:uv}

\subsection{The ultraviolet row}

A usable ultraviolet certificate must state at least:
\begin{enumerate}[label=(\arabic*)]
  \item the bare or regulated objects;
  \item the renormalization or continuum-limit prescription;
  \item the symmetry or Ward identities maintained or recovered;
  \item the norm and observable class in which convergence is claimed;
  \item the energy range and truncation error; and
  \item whether the construction is an effective theory, a perturbative
  completion, or a nonperturbative completion.
\end{enumerate}

Wilsonian coarse-graining and effective Lagrangians
\cite{WilsonKogut1974,Polchinski1984}, an asymptotic fixed point
\cite{Reuter1998}, modular finiteness, a discrete continuum limit, and an entire
spectral form factor are not the same certificate. They may share low-energy
predictions without sharing ultraviolet operators.

\begin{proposition}[Low-energy agreement does not select spectral damping]
\label{prop:uv-nonuniqueness}
Fix \(m>0\), an accessible momentum bound \(E>0\), and \(\Lambda>0\). Define
\[
 G_0(p)=\frac{1}{p^2+m^2},
 \qquad
 G_\Lambda(p)=\frac{e^{-p^2/\Lambda^2}}{p^2+m^2}.
\]
Then for \(|p|\leq E\),
\[
 |G_0(p)-G_\Lambda(p)|
 \leq \frac{E^2}{\Lambda^2m^2}.
\]
Consequently, low-energy agreement to any fixed tolerance does not determine
whether the ultraviolet kernel contains smooth exponential damping.
\end{proposition}

\begin{proof}
Use \(0\leq 1-e^{-x}\leq x\) for \(x\geq0\):
\[
 |G_0-G_\Lambda|
 =
 \frac{1-e^{-p^2/\Lambda^2}}{p^2+m^2}
 \leq
 \frac{p^2}{\Lambda^2(p^2+m^2)}
 \leq
 \frac{E^2}{\Lambda^2m^2}.
\]
\end{proof}

The example does not claim that both kernels define the same complete theory.
It proves the narrower logical point: finite predictivity on a bounded range
cannot force one ultraviolet mechanism. It also shows why a four-dimensional
smooth UV filter cannot be inferred merely from suppression of an internal
operator sector.

\section{When events are actually discrete}
\label{sec:events}

An outcome label, a quantum jump, a first-passage time, a detector click, and an
element of a causal set are different objects. Their discreteness arises from
different assumptions.

\begin{theorem}[Transverse crossings are finite on compact time intervals]
\label{thm:transverse-events}
Let \(f\in C^1([0,T],\R)\). If
\[
 f(t)=0\quad\Longrightarrow\quad f'(t)\neq0,
\]
then the zero set of \(f\) in \([0,T]\) is finite.
\end{theorem}

\begin{proof}
Every zero is isolated by the inverse-function theorem. If there were
infinitely many zeros in the compact interval, they would possess an
accumulation point \(t_\ast\). Continuity gives \(f(t_\ast)=0\). But
\(f'(t_\ast)\neq0\) makes \(t_\ast\) an isolated zero, contradicting
accumulation.
\end{proof}

This theorem supports an event interpretation for smooth, transverse
first-exit models on finite horizons. Without the hypothesis, zeros may
accumulate; stochastic paths may revisit a boundary infinitely often; reset
dynamics may exhibit Zeno behavior; and an outcome instrument may be discrete
without any underlying boundary crossing.

A causal set requires more still: a partial order and local finiteness of order
intervals. Those are defining structural data in the causal-set program
\cite{BombelliEtAl1987}. A countable list of basin transitions does not by
itself supply either property. Conversely, discrete spacetime proposals cannot
be rejected merely by saying that a lattice violates Lorentz symmetry; the
actual construction and continuum approximation must be examined.

\section{Gauge and matter claims require global data}
\label{sec:matter}

\subsection{The gauge-and-matter row}

At minimum, a matter certificate must specify:
\[
 \axis{G}
 =
 \bigl(
 G_{\rm global},\{\rho_i\},\mathcal H_{\rm ferm},
 \mathfrak s,\mathcal L,\mathcal A_{\rm tot},
 \mathcal O_{\rm allowed},\nu
 \bigr).
\]
Here \(G_{\rm global}\) is the global gauge group rather than only its Lie
algebra; \(\rho_i\) are representations; \(\mathfrak s\) records spin or
Spin\(^{c}\) data; \(\mathcal L\) records bundles and characteristic classes;
\(\mathcal A_{\rm tot}\) is the complete anomaly including inflow and
counterterms; \(\mathcal O_{\rm allowed}\) is the operator-selection rule; and
\(\nu\) fixes charge normalization.

\subsection{What a compact circle group does prove}

\begin{proposition}[Weights of a specified compact circle]
\label{prop:u1-weights}
Every finite-dimensional continuous unitary representation of \(U(1)\) is a
direct sum of characters
\[
 z\longmapsto z^n,
 \qquad n\in\mathbb Z.
\]
Thus charges are integer weights after a generator and normalization have been
fixed.
\end{proposition}

\begin{proof}
The commuting unitary matrices in the image can be simultaneously
diagonalized. Each one-dimensional continuous character of \(U(1)\) has the
form \(e^{i\theta}\mapsto e^{in\theta}\) with \(n\in\mathbb Z\).
\end{proof}

This proposition does not derive the observed hypercharges. Fractional
numerical values depend on normalization and on the global quotient connecting
the \(U(1)\) factor to nonabelian centers. Replacing \(U(1)\) by the additive
group \(\R\) also removes the integer-character conclusion. The common-circle
idea in MTT is therefore potentially useful only after the same physical line,
connection, and global group action are identified.

\subsection{Anomalies and operator selection}

Gauge consistency constrains the \emph{total} anomaly. It need not require the
bare fermion polynomial to vanish term by term: Green--Schwarz inflow is a
standard counterexample \cite{GreenSchwarz1984}. A theorem must state the
dimension, chiral content, regularization, local and global anomalies, inflow,
and boundary conditions.

\begin{proposition}[Conditional bundle obstruction]
\label{prop:bundle-obstruction}
Suppose fields \(\psi_i\) are sections of associated bundles \(E_i\), and a
candidate local monomial with constant scalar coefficient transforms as a
section of a product bundle \(E_{\mathcal O}\). If \(E_{\mathcal O}\) admits no
gauge-invariant identification with the scalar-density bundle, that monomial
cannot define a global gauge-invariant action term with that coefficient.
Conversely, such an identification removes this obstruction but does not force
the term to occur or fix its coefficient.
\end{proposition}

\begin{proof}
A global action density must be a gauge-invariant scalar density. Without an
equivariant identification of \(E_{\mathcal O}\) with that target, local
expressions fail to glue to a global scalar. If an identification exists, the
gluing obstruction is absent, but dynamics and additional symmetries still
decide whether the coefficient vanishes.
\end{proof}

This is the correct scope for topology-only exclusion. It cannot be promoted to
the statement that all dangerous baryon- or lepton-number violating operators
are absent in every admissible theory. Such operators are systematically
present in effective-field-theory classifications unless additional symmetries
or bundle rules exclude them \cite{Weinberg1979}. Nor do anomaly freedom,
chirality, and charge quantization alone prove that the Standard Model is the
unique or near-minimal solution.

\section{The mathematics of a multi-axis classification}
\label{sec:classification}

\subsection{Exact classes}

Let \(\mathfrak T\) be a set of typed model records. For each axis
\(k\in K\), let
\[
 I_k:\mathfrak T\to Q_k
\]
be the invariant extracted from that row, and let \(\simeq_k\) be an
equivalence relation on \(Q_k\).

\begin{definition}[Axis-set equivalence]
For \(S\subseteq K\), define
\[
 a\sim_S b
 \quad\Longleftrightarrow\quad
 I_k(a)\simeq_k I_k(b)
 \ \text{for every }k\in S.
\]
\end{definition}

\begin{theorem}[Intersection and refinement of classification axes]
\label{thm:axis-refinement}
For every \(S\subseteq K\), the relation \(\sim_S\) is an equivalence
relation. If \(S\subseteq T\subseteq K\), then every \(\sim_T\)-class is
contained in a \(\sim_S\)-class.
\end{theorem}

\begin{proof}
The relation \(\sim_S\) is the intersection of the equivalence relations
induced by the maps \(I_k\), so it is reflexive, symmetric, and transitive.
If \(a\sim_T b\), then the defining conditions hold for every \(k\in T\),
hence for every \(k\in S\).
\end{proof}

The theorem is elementary, but it carries the main methodological message:
adding measurement, ultraviolet, matter, or validation rows can only refine a
class. Agreement on one row cannot establish agreement on a stronger set.

\begin{proposition}[No inference between independent axes]
\label{prop:no-cross-axis}
Let \(Q_1,\ldots,Q_n\) each contain at least two elements, and take
\(\mathfrak T=Q_1\times\cdots\times Q_n\) with \(I_k\) the coordinate
projection. For any proper subset \(S\subsetneq\{1,\ldots,n\}\), there exist
\(a,b\in\mathfrak T\) such that \(a\sim_S b\) but \(I_j(a)\neq I_j(b)\) for
some \(j\notin S\).
\end{proposition}

\begin{proof}
Choose \(a\) and \(b\) equal in every coordinate in \(S\), and choose distinct
values in one coordinate \(j\notin S\).
\end{proof}

Physical axes may be coupled by nontrivial theorems. The proposition says that
the coupling must be proved; it is not supplied by the vocabulary.

\subsection{Approximate classes}

A metric tolerance such as
\[
 d_k\bigl(I_k(a),I_k(b)\bigr)\leq\varepsilon_k
\]
is generally not transitive. It defines a comparison neighborhood or graph,
not automatically an equivalence class. To speak of a genuine universality
class one needs an exact equivalence, a quotient construction, a common limit,
or a controlled transitive closure whose accumulated error remains acceptable.
The local coherent-sector comparison paper proves the relevant resolvent bound
and explicitly tracks chainwise error accumulation
\cite{NeroLocalRobustness2026}. It does not license unlimited chains.

\begin{definition}[Conditional coherent comparison class]
A family \(\mathfrak C\subseteq\mathfrak T\) is a \emph{conditional coherent
comparison class} on axis set \(S\) when:
\begin{enumerate}[label=(\roman*)]
  \item every model has a complete witness on \(S\);
  \item the transport maps and comparison relations are fixed before testing;
  \item exact equivalence or a uniformly controlled common-limit theorem is
  proved on each axis in \(S\); and
  \item the validation row states what physical agreement follows.
\end{enumerate}
\end{definition}

This definition deliberately avoids the word ``inevitable.'' Membership is a
verified property of supplied records.

\section{What the major external frameworks actually contribute}
\label{sec:external}

The comparison below records native structures, not verdicts on entire research
programs.

\begin{table}[htbp]
\centering
\scriptsize
\renewcommand{\arraystretch}{1.18}
\begin{tabularx}{\textwidth}{@{}P{0.17\textwidth}Y Y@{}}
\toprule
Framework & Native contribution relevant here & What does not follow automatically \\
\midrule
Operational quantum theory &
Instruments, completely positive operations, outcome probabilities, and
sequential measurement \cite{DaviesLewis1970}. &
A unique ontology, objective event source, or irreversible microscopic law. \\

Objective collapse &
Explicit stochastic modification producing localization and macroscopic
classicality \cite{GRW1986}. &
Derivation from generic record stability or equivalence to coarse-graining. \\

Monitored circuits &
Model-specific entanglement phases and critical behavior under measurement
\cite{SkinnerRuhmanNahum2019}. &
A universal solution of measurement or one class for all instruments. \\

EFT and RG &
Controlled low-energy expansions, coarse-graining, running couplings, and
universality near fixed points \cite{WilsonKogut1974,Polchinski1984}. &
A unique UV completion, literal irreversibility of every beta-function ODE, or
smooth spectral damping. \\

Asymptotic safety &
A candidate non-Gaussian fixed-point mechanism formulated through
scale-dependent effective actions \cite{Reuter1998}. &
Equality to string modular control or an entire form-factor kernel. \\

Loop quantum gravity &
Background-independent quantum geometry and a mathematically specific
kinematical representation \cite{AshtekarLewandowski2004}. &
An MTT projector, measurement completion, or full low-energy equivalence without
an explicit bridge. \\

Causal sets &
A locally finite partial order proposed as microscopic spacetime structure
\cite{BombelliEtAl1987}. &
Derivation from generic basin crossings or failure of Lorentzian approximation. \\

Noncommutative geometry &
A spectral triple and spectral action once the algebra, Hilbert space, and
Dirac operator are supplied \cite{ChamseddineConnes1997}. &
Selection of that triple from a generic coherent sector or uniqueness of the
observed Standard Model. \\
\bottomrule
\end{tabularx}
\caption{A framework should be compared through its own typed objects.}
\label{tab:frameworks}
\end{table}

The table reverses the old ``partial shadows'' rhetoric. Similarity is now a
question to be certified, not an interpretation imposed in advance.

\section{The present MTT position}
\label{sec:mtt}

MTT currently supplies several distinct layers:
\begin{enumerate}[label=(\arabic*)]
  \item The Foundations paper defines the closure/admissibility program and its
  typed status boundaries \cite{NeroFoundation2026}.
  \item The Projection--Admissibility paper distinguishes descent,
  recoverability, and structural constraints on an effective description
  \cite{NeroProjection2026}.
  \item The coherent-sector reduction paper proves a domain-explicit
  single-model Feshbach and projector certificate, including the special case
  in which an exact spectral projector makes the off-diagonal correction vanish
  \cite{NeroCoherentReduction2026}.
  \item The local robustness paper transports retained resolvents to one
  reference Hilbert space and proves bounded local comparison, while separating
  projector, operator, dynamical, observable, and physical equivalence
  \cite{NeroLocalRobustness2026}.
  \item The measurement paper separates physical disturbance, outcome
  completion, record stabilization, and the still-open source of a general
  outcome law \cite{NeroMeasurement2026}.
\end{enumerate}

These are meaningful advances, but they do not close all seven witness rows for
MTT or for any external framework. In particular:
\begin{itemize}
  \item a selected local projector is not a universal measurement instrument;
  \item an internal spectral filter is not by itself a four-dimensional UV
  completion;
  \item a rank or bundle flag is not yet the observed matter spectrum;
  \item a finite event encoding is not automatically a causal set;
  \item local resolvent closeness is not physical equivalence; and
  \item one successful MTT realization would not prove that all successful
  theories had to use the same microscopic mechanism.
\end{itemize}

The phrase \emph{coherent universality} may still be used, but only with an axis
set and certificate. For example, ``local norm-resolvent coherent universality
on \(\Omega\)'' is precise. ``All viable physics belongs to one coherent
universality class'' is not presently a theorem.

\section{A practical comparison protocol}
\label{sec:protocol}

For a proposed comparison between MTT and another framework:
\begin{enumerate}[label=\textbf{Step \arabic*.},leftmargin=4.8em]
  \item Freeze the exact source versions and their provenance hashes.
  \item Choose the axis set \(S\). Do not add untested axes in the conclusion.
  \item Fill both witness records, including domains, norms, units, and time or
  energy ranges.
  \item Construct explicit transports to common objects.
  \item Prove commutation, convergence, or a quantitative error bound.
  \item Map preparations, states, observables, and predictions.
  \item Reserve at least one independent or held-out test.
  \item Record failed rows and breakdown surfaces as part of the result.
\end{enumerate}

\subsection{Minimum machine-readable payload}

A computational record should include:
\begin{verbatim}
model_id, source_hashes, axis_set,
state_space, domains, transports, comparison_norm,
parameter_region, time_or_energy_horizon,
error_bounds, observable_map, unit_map,
status, failed_rows, verifier, artifact_hashes
\end{verbatim}
The payload does not replace the proof. It prevents a later comparison from
silently changing its objects or standard.

\section{Completion and falsifiability}
\label{sec:falsifiability}

The framework can fail in informative ways.
\begin{enumerate}[label=(F\arabic*)]
  \item A claimed descent map is injective, while the argument requires
  information loss.
  \item The recovery defect vanishes in the declared physical decoder class.
  \item Record regions have zero margin or fail finite-horizon retention.
  \item The ultraviolet row lacks a continuum or error certificate.
  \item event times accumulate, or no causal order is supplied.
  \item the global gauge group, anomaly total, or bundle product contradicts
  the proposed matter claim.
  \item transported observables disagree beyond the declared uncertainty.
\end{enumerate}

Conversely, completing all rows would justify a strong result: not that two
theories are ontologically identical, but that they are equivalent on a
declared physical interface to a certified accuracy.

\section{Conclusion}
\label{sec:conclusion}

The strongest correction in this paper is also the simplest. Empirical demands
constrain theories, but they do not name the mathematical mechanism satisfying
them. Irreversible records do not alone force a noninjective projection. Stable
records do not uniquely force attractor basins or sharp knees. Finite
predictivity does not select smooth spectral damping. Event-like outcomes do
not automatically define a causal set. Gauge consistency and compact-circle
weights do not uniquely determine the observed matter sector.

What survives is a rigorous and reusable classification method. The seven-axis
witness record shows exactly what is common, what has merely been translated,
what is quantitatively controlled, and what remains open. The exact
classification theorem says that adding axes refines rather than preserves a
class; the counterexamples show why cross-axis conclusions require their own
bridges.

For MTT, this is not a retreat from comparison. It is what makes comparison
credible. The released operator results now occupy their correct local layer,
while measurement, ultraviolet completion, event structure, matter selection,
and empirical equivalence remain independently testable rows. A future theorem
may couple some of those rows. Until it does, the correct claim is conditional
classification, not inevitability.

\bibliographystyle{plain}
\bibliography{main}

\end{document}
